\section{相关研究现状}

随着LLM的攻击影响由文本输出扩展到真实工具动作及跨步骤状态传播，防御研究的控制对象也相应发生变化：早期方法主要在输入侧区分或净化不可信内容，随后通过训练提高模型对指令边界的稳定识别能力，进一步发展为在Agent运行时直接约束工具动作、参数来源与执行状态。以下首先梳理间接提示注入攻击的演化过程，再按照“提示工程与输入净化—训练与安全对齐—Agent系统架构设计”的递进关系分析相关防御方法研究现状。

\subsection{检索增强生成技术研究现状}

提示注入研究首先关注攻击者如何通过不可信输入改变模型对原始指令的解释与执行\cite{liu2024formalizing}；越狱与指令劫持研究进一步表明，模型可能在角色诱导、对抗性指令或上下文冲突下偏离原有安全策略\cite{shen2024do,zeng2024howjohnny2}。与主要作用于当前模型响应的直接注入或越狱攻击不同，Greshake等\cite{greshake2023notwhat}提出的间接提示注入攻击将恶意指令嵌入网页、邮件和文档等外部数据源，使Agent在完成正常任务、读取工具返回值时将数据中的内容解释为控制指令，并进一步影响后续工具调用。由此，攻击目标从诱导模型生成特定文本，扩展为操纵Agent的动作选择、关键参数和外部副作用。此后，相关攻击主要沿以下方向发展：

\textbf{（1）工具选择与参数操纵。}攻击者可以诱导Agent选择攻击者指定的工具\cite{shi2026pitoolselection}，也可以在工具名称符合任务表面要求的情况下操纵关键参数或利用已授予权限实施非预期动作\cite{zhang2026evaluatingprivilege}。因此，仅检查工具名称或静态权限集合不足以判断实际调用是否符合用户任务。

\textbf{（2）多步骤与工作流组合。}Kothamasu等\cite{kothamasu2026hidden}将恶意意图拆分为多个表面上相互独立的子任务，并分散到不同工具返回内容中，使Agent在多步规划过程中将其组合为完整攻击链。该类攻击不要求单个观测或单次调用独立呈现完整恶意目标，而是利用跨步骤的信息传递与状态积累实现攻击效果。

\textbf{（3）安全信息侦察与持久化操纵。}攻击者可利用对抗性提示诱导Agent泄露系统提示、上下文记忆或内部策略，以获取后续攻击所需的任务边界信息\cite{hui2024scp}；也可向长期记忆、检索知识库或训练数据植入投毒内容和后门触发器，使恶意影响在后续任务中持续存在\cite{chen2024agentpoison,bi2024backdoor}。

\textbf{（4）面向防御机制的自适应攻击。}在获知系统是否拒绝、如何反馈或如何净化输入后，攻击者可以针对防御边界迭代调整注入载荷。相关评测表明，自适应攻击会改变部分静态防御在原始攻击下呈现的安全效果\cite{zhan2025adaptive}，因而防御机制还需要考虑拒绝反馈、重规划与再次调用所形成的攻击面。

% 在评测层面，ToolEmu利用大模型模拟沙箱评估Agent工具使用风险\cite{ruan2024toolemu}；InjecAgent\cite{zhang2024injectagent}、AgentDojo\cite{balunovic2024agentdojo}、Agent Security Bench\cite{zhang2025asb}和AgentDyn\cite{li2026agentdyn}则面向间接提示注入攻击，构造包含外部不可信数据、工具调用和任务效用判定的专门化评测环境。

上述演化表明，Agent工具调用链中的间接提示注入已由单轮文本冲突发展为动作与参数操纵、跨步骤传播和持久状态影响并存的运行时安全问题。其防御既需要在副作用发生前判断当前动作是否具有可信授权，也需要考虑当多步隐蔽攻击的影响已进入运行轨迹时，及时识别其传播证据并予以阻断或干预，以控制后续损失并恢复Agent系统运行。

\subsection{检索增强生成系统投毒攻击研究现状}

针对外部数据可能被模型解释为控制指令的问题，最先形成的是输入侧防御。此类方法主要通过调整提示文本组织方式或净化外部内容，引导模型区分可信指令与不可信数据，通常不改变模型参数和Agent执行架构，具有较低的接入成本。Schulhoff\cite{schulhoff2024sandwich}提出Prompt Sandwiching，在不可信内容前后重复插入用户原始指令，使模型在读取外部数据后重新关注初始任务。Hines等\cite{hines2024defending}提出Spotlighting，将单纯重复指令推进为使用特殊分隔符、字符编码或数据标记显式区分指令与外部内容。随后，Shi等\cite{shi2025promptarmor}提出PromptArmor，进一步使用额外的大模型检测并移除外部输入中的注入指令，使防御由提示结构调整扩展到语义净化。Agent Security Bench\cite{zhang2025asb}则将分隔符、提示夹心、指令预防、复述和乱序等方法纳入统一基线，推动了输入侧方法在相同攻击与任务条件下的系统比较。

上述方法\cite{schulhoff2024sandwich,hines2024defending,zhang2025asb,shi2025promptarmor}主要依赖大语言模型对分隔符、重复指令和净化结果的概率性遵从，未在真实工具执行边界形成独立的动作约束。已有评测\cite{zhan2025adaptive,balunovic2024agentdojo,zhang2024injectagent}表明，其防御效果会随任务、攻击表达方式和自适应载荷发生变化。因此，输入侧方法虽然降低了部署门槛，但也暴露出仅依靠推理时提示信号难以稳定维持指令边界的问题，促使后续研究将防御能力进一步固化到模型训练与专用检测器中。

\subsection{检索增强生成系统防御研究现状}

为克服提示信号易受模型概率性行为影响的问题，研究重点随后转向通过训练或微调，使指令层级和注入识别能力成为较稳定的模型行为。在大模型安全对齐方面，Wallace等\cite{wallace2024instruction}提出Instruction Hierarchy，显式定义系统消息、用户输入与第三方工具或数据输出的优先级，并构造训练数据使模型在指令冲突时优先遵循高权限指令。Chen等\cite{chen2025struq}提出StruQ，将提示层的分隔思想推进为结构化通道和特殊标记，并训练模型识别相应边界。SecAlign\cite{chen2025secalign}和Meta SecAlign\cite{chen2025metasecalign}进一步将提示注入防御建模为偏好优化问题，使模型在包含注入指令的输入下偏好安全响应。随着防御对象由文本响应延伸到Agent动作，Mou等\cite{mou2026toolsafe}提出ToolSafe，通过多任务防御反馈强化学习训练步骤级安全护栏，在每次工具调用前进行风险判定。

与直接改造主模型并行发展的另一条路线是训练独立检测器。ProtectAI\cite{protectai2024deberta}基于DeBERTa-v3-base微调二分类模型，以较低推理开销识别注入文本；在此基础上，PIGuard\cite{li2025piguard}进一步关注高检出率伴随的过度防御问题，以兼顾攻击检出率与良性输入误报率；DataSentinel\cite{liu2025datasentinel}则将研究重点推进到适应性攻击，采用极小极大训练框架建模攻击者与检测器之间的对抗过程。

相较于提示工程，训练类方法将部分防御能力从临时提示固化为模型或检测器的统计行为，但其效果仍受训练数据分布和模型更新成本影响。步骤级护栏\cite{mou2026toolsafe}主要输出安全判定，独立检测器\cite{li2025piguard,liu2025datasentinel}则主要处理文本片段；二者均需要与Agent的最终动作、参数来源和真实执行入口进一步绑定，才能将检测结果转化为强制性的运行时约束。由此，研究重心继续从“模型能否识别风险”转向“系统能否在风险识别后约束实际执行”，推动了Agent系统架构层防御的发展。


综上所述，现有LLM Agent间接提示注入防御已由提示层的概率性引导逐步发展到训练型检测和系统级执行约束。围绕调用前防御，已有研究覆盖了任务相关性检查、符号权限策略、参数来源契约、隔离执行和反事实归因等机制，为进一步区分证据等级和契约判据提供了基础。围绕跨步骤响应，已有研究能够利用执行图或反事实分支识别外部内容对动作的影响，但从可核验运行证据出发完成路径定位、可执行阻断、受影响状态恢复和处置后验证仍需要形成一致的运行闭环；而在确定性规则无法唯一消歧的候选影响边上，如何利用图结构与语义信息自动排序可疑程度并引导审查资源分配，尚有待进一步研究。基于上述研究脉络，本研究分别从基于执行轨迹图分级证据判据的在线防御和基于图-语义联合表示的攻击定位与溯源两个方向展开研究。

